\section{Supplementary proofs and scope ledger}
\label{app:proofs}

\subsection{Primitive-root regrouping with weights}

Let $F=\sum_{e\in\E_A}w(e)$.  The coefficient of $z^n$ in
$\sum_{n\ge1}z^nF^n/n$ is a sum over based words of length $n$.  Every based
word has a unique primitive cyclic root $\gamma$ of length $\ell$ and a
unique repetition number $r=n/\ell$.  A primitive necklace of length $\ell$
has exactly $\ell$ based rotations.  Therefore its total coefficient in the
based-word sum is
\[
  \frac{\ell}{r\ell}w(\gamma)^r=\frac{w(\gamma)^r}{r}.
  \tag{A.1}
\]
This proves (3.10) as a formal identity.  In the analytic specialization it
also holds wherever the trace-log converges absolutely.

The formula makes the scalar-sign firewall unavoidable.  If a primitive
weight is negative, its $r$th term in (A.1) contains the $r$th power of that
minus sign.  A fixed chain parity is a different trace convention.

\subsection{Complete \texorpdfstring{$p^2q^2$}{p2q2} necklace census}

The content equation for a word with counts $(n_a,n_b,n_c)$ is
\[
  n_a+n_c=2,\qquad n_b+n_c=2.
  \tag{A.2}
\]
Thus $(n_a,n_b,n_c)$ is one of $(2,2,0)$, $(1,1,1)$, or $(0,0,2)$.

For $(2,2,0)$, the binary necklaces are represented by $aabb$ and $abab$.
The latter has least period two.  For $(1,1,1)$, cyclic orders modulo rotation
are represented by $abc$ and $acb$, both of least period three.  For
$(0,0,2)$, the only word is $cc$, of least period one after taking the
primitive root.  This proves that (4.3) and the two $r=2$ terms exhaust the
coefficient.

\subsection{Irreducible decomposition of the \texorpdfstring{$S_3$}{S3} residual}

The irreducible character table in class order
$1,(12),(123)$ is
\[
\begin{array}{c|rrr}
 &1&(12)&(123)\\ \hline
\one&1&1&1\\
\sgnrep&1&-1&1\\
\stdrep&2&0&-1
\end{array}
\tag{A.3}
\]
and the virtual residual is $(0,0,3)$.  Its inner products with the three
irreducible characters, using class sizes $1,3,2$, are
\[
  \langle\chi,\one\rangle=1,
  \qquad
  \langle\chi,\sgnrep\rangle=1,
  \qquad
  \langle\chi,\stdrep\rangle=-1.
  \tag{A.4}
\]
Hence $\chi=\one\oplus\sgnrep-\stdrep$.  Its dimension is zero, explaining
why the scalar determinant cannot see it.

\subsection{Koszul and Harrison index conventions}

The Koszul computation in \Cref{thm:tor} concerns Tor of the augmentation
module.  HKR in (6.4) concerns Hochschild homology with coefficients in the
algebra.  Harrison homology and André--Quillen homology use another quotient
and an index shift: in characteristic zero, higher Harrison groups encode
the derived indecomposable commutative direction.  For a polynomial algebra,
the cotangent complex is the projective module of K\"ahler differentials in
degree zero, so higher derived functors vanish.

These statements are compatible, not contradictory.  They answer different
questions.  The mixed Tor and Hochschild exterior classes prove that an
atom-only commutative quotient is not a quasi-isomorphic replacement for the
full bar/Hochschild recurrence ledger.

\subsection{Trace-class extension of the acyclic lemma}

\Cref{thm:acyclic} was stated in finite dimension, which is all the exact
control requires.  The same graded-commutator proof extends to a countable
complex only when the products in (6.9) are trace class and supertrace
cyclicity is justified.  No such infinite chain complex is source-locked for
SD-C17, so we do not claim the extension here.

\subsection{Scope and nonclaims}

The paper proves no statement about:

\begin{itemize}
\item zeros of $\zeta$ outside the Euler-product half-plane;
\item the completed function $\xi$, its functional equation, or Gamma factor;
\item a Riemann--von Mangoldt counting law or Weil compression;
\item a self-adjoint, unitary, scattering, or Hamiltonian operator;
\item all possible representation-valued symbolic determinants;
\item all chain-enhanced grammars that add new source data before cyclic
  reduction.
\end{itemize}

The countable identity $D_\infty=1/\zeta$ is valid only in $\Ree s>1$ under
the absolute alphabet sum (3.7).  Analytic continuation of the already-known
Euler product is not credited to the symbolic shift.

The route tuple remains
\[
\begin{split}
(&\texttt{A0\_ANALYTIC\_ARITHMETIC\_ORIGIN},
  \texttt{A1\_FAIL},\\
 &\texttt{A2\_ANALYTIC\_DETERMINANT},
  \texttt{A3\_FAIL},\texttt{A4\_FAIL}),
\end{split}
\]
with \texttt{ROUTE\_A\_REJECTED} and Route B locked.
